\documentclass[11pt]{article}
\usepackage[margin=1.05in]{geometry}
\usepackage{amsmath,amsthm,amssymb,mathtools,bm}
\usepackage{booktabs,array,enumitem,microtype}
\usepackage[colorlinks=true,linkcolor=blue!60!black,citecolor=blue!60!black,urlcolor=blue!60!black]{hyperref}
\theoremstyle{plain}
\newtheorem{theorem}{Theorem}[section]\newtheorem{proposition}[theorem]{Proposition}\newtheorem{lemma}[theorem]{Lemma}
\newtheorem{corollary}[theorem]{Corollary}
\theoremstyle{definition}\newtheorem{definition}[theorem]{Definition}\newtheorem{remark}[theorem]{Remark}\newtheorem{example}[theorem]{Example}
\newcommand{\R}{\mathbb R}\newcommand{\E}{\mathbb E}\newcommand{\Var}{\operatorname{Var}}\newcommand{\Cov}{\operatorname{Cov}}
\newcommand{\relu}{\operatorname{relu}}\newcommand{\diag}{\operatorname{diag}}\newcommand{\tr}{\operatorname{tr}}
\newcommand{\kap}{\kappa}\newcommand{\Sym}{\operatorname{Sym}}\newcommand{\Hpl}{\mathrm{He}}
\title{\textsc{Localization in the commutant, stage 9:}\\ \large the mean of a random \textsc{ReLU} network as a central localization,\\ with the exact slice transport, the collective-mode sensitivity law\\ and the capacity of input tilts}
\author{Stage 9 of the quantum-cut / noncommutative-geometry programme}
\date{9 October 2026}
\begin{document}\maketitle
\begin{abstract}
We apply the stage-8 principle (a localization process is a positive martingale in the commutant, whose central part is
the transverse measure on the leaf space) to the WhestBench problem: the per-neuron mean $\E\,h_L$ of a bias-free
He-initialised \textsc{ReLU} multilayer perceptron with standard Gaussian input, $n=1024$, $L=16$. The cumulant chains that
dominate the leaderboard are read as localizations: the births of non-Gaussianity at the \textsc{ReLU} kinks are the
linear-tilt responses of the gate algebra, the chain is the trickle-down expansion of that algebra's product state, and
the one object that no first-order chain carries correctly is the \emph{central} variable---the collective activity
mode, a scale mixture---which is exactly the leaf-space part of the localization. Main results. (A) Input-side linear
tilts have localization capacity $O(k/n)$ for the born non-Gaussianity and cost $\ge 2k+1$ chain evaluations: input
localization is useless, which forces the neuron-side (gate) decomposition. (B) A sensitivity law: the output mean
responds to a relative error $\varepsilon$ of the collective mode's variance with error $c\,\varepsilon$ per layer,
$c\approx0.3$, so the collective mode must be accurate to $10^{-3}$ per layer while the bulk may be off by $10^{-1}$;
on the official networks the chain's covariance is accurate everywhere except in that one direction. (C) The exact
first-order transport of the $(2,1)$ slices and of the diagonal through \textsc{ReLU}, in closed form, including the
transposed-slice and the $\kap_4\!\to\!\kap_3$ feeds, and the exact-by-homogeneity transport of the scale mixture. (D)
On the official networks, the terms prescribed by the central/Cartan split account for the measured error ladder
$8.5\cdot10^{-7}\to5.7\cdot10^{-7}$ of a clean chain and locate the remaining deficit in the transported third-cumulant
slices (25\% low at depth). Every identity is checked by the accompanying script.
\end{abstract}
\tableofcontents
\section{The problem and the principle}
\subsection{The estimation problem}
Fix weights $W_1,\dots,W_L\in\R^{n\times n}$, drawn once from $N(0,2/n)$ (He initialisation) and then held fixed. For an
input $x\in\R^n$ put $h_0=x$, $z_l=W_l h_{l-1}$, $h_l=\relu(z_l)$. The target is the vector of post-activation means
\[
u\;=\;\E_{x\sim N(0,I_n)}\,h_L(x)\in\R^n ,
\]
scored by the mean squared error over the $n$ neurons, multiplied by $\max(0.1,\,C/B)$ with $C$ the metered floating
point count and $B=2^{41}$. The map $F=h_L$ is positively homogeneous of degree one and piecewise linear. The reference
answer is a Monte Carlo average over $10^9$ inputs, whose per-neuron noise ($\sim10^{-5}$) is below every error
discussed here.

\subsection{Estimators are functions on the state space; exact ones are affine}
Let $\mathcal S$ be the set of probability laws $\nu$ on $\R^n$ with enough moments. Every deterministic estimator we
consider is a map $E:\mathcal S_0\to\R^n$ defined on a subfamily $\mathcal S_0\subseteq\mathcal S$ that contains the point
masses $\delta_x$ and the Gaussians $N(m,\Sigma)$; a cumulant chain is defined on Gaussians by running its recursion
from the given $(m,\Sigma)$, and on point masses by the forward pass, $E(\delta_x)=F(x)$. The exact answer
\[
u(\nu)=\nu(F)=\int F\,d\nu
\]
is \emph{affine} on $\mathcal S$. This is the whole content of exactness: in the commutant language of \cite{stage8},
$u$ is the pairing of the state $\nu$ with the observable $F$, and a localization process of $\nu$ (a positive
martingale of Radon--Nikodym derivatives in the commutant) turns $u$ into a martingale. A chain is not affine: it is
exact on point masses and wrong on mixtures, and the gap between the two is what we have to compute.

\begin{definition}[Localization process]
A localization process of $\nu$ is a sequence or continuous family $(\nu_t)_{t\ge0}$ of random laws with $\nu_0=\nu$,
adapted to a filtration $(\mathcal F_t)$, such that $\E[\nu_{t+s}(g)\mid\mathcal F_t]=\nu_t(g)$ for every bounded $g$,
and \emph{complete} if $\nu_t\to\delta_X$ almost surely for a random $X\sim\nu$ \cite{CE,stage8}. Three schemes are used
below: Eldan's isotropic tilt, $\nu_t=\mathrm{Law}(x\mid y_t)$ with $y_t=tX+B_t$, so that for $\nu=N(0,I)$ one has
$\nu_t=N(m_t,\,I/(1+t))$ with $m_t$ a martingale of quadratic variation $d\langle m\rangle_t=(1+t)^{-2}I\,dt$; the
\emph{directional} tilt, the same observation with $y_t=t\,\langle v,X\rangle+B_t$ for a unit vector $v$, which leaves
$\nu_t=N(m_t,\,I-\tfrac{t}{1+t}vv^{\top})$; and pinning, which reveals the coordinates of $X$ one at a time.
\end{definition}

\begin{theorem}[Defect--path formula]\label{thm:defect}
Let $E:\mathcal S_0\to\R^n$ be twice continuously Fr\'echet differentiable along the Gaussian family and exact on point
masses, and let $(\nu_t)$ be a complete localization process of $\nu\in\mathcal S_0$ whose components stay in
$\mathcal S_0$ and have bounded quadratic variation. Then
\begin{equation}\label{eq:defect-path}
u(\nu)-E(\nu)\;=\;-\,\E\int_0^\infty \tfrac12\,E''(\nu_t)\big[d\langle\nu\rangle_t\big]
\;=\;-\,\E\int_0^\infty\mathrm{drift}_t\big(E(\nu_t)\big)\,dt .
\end{equation}
In words: the error of the estimator at $\nu$ is minus the expected total drift of the estimator along any complete
localization of $\nu$, and that drift is the second variation of $E$ in the directions the scheme localizes.
\end{theorem}

\begin{proof}
$u(\nu_t)$ is a martingale (the law of iterated expectations, or affinity of $u$ plus the martingale property of
$\nu_t$), so $u(\nu)=\E\,u(\nu_\infty)=\E F(X)$. By It\^o's formula for the semimartingale $\nu_t$ in the parametrisation
of $\mathcal S_0$, $dE(\nu_t)=E'(\nu_t)[d\nu_t]+\tfrac12E''(\nu_t)[d\langle\nu\rangle_t]$, and the first term has zero
mean. Hence $\E\,E(\nu_\infty)-E(\nu)=\E\int_0^\infty\tfrac12E''(\nu_t)[d\langle\nu\rangle_t]$. Exactness on point
masses gives $E(\nu_\infty)=F(X)$ almost surely, so $\E\,E(\nu_\infty)=u(\nu)$. For a discrete scheme the same identity
holds with the drift replaced by the conditional Jensen gaps $\E[E(\nu_{t+1})\mid\mathcal F_t]-E(\nu_t)$.
\end{proof}

\begin{corollary}[Eldan's path for a homogeneous network]\label{cor:eldan}
Write $e(y)=E(N(y,I))$ for a chain $E$ that is exact on point masses and transforms like $F$ under scaling,
$E(N(m,sI))=\sqrt s\,e(m/\sqrt s)$. Then with $\delta(y)=e(y)-y\cdot\nabla e(y)-\Delta e(y)$,
\begin{equation}\label{eq:eldan}
u-e(0)\;=\;-\tfrac12\int_0^\infty(1+\tau)^{-3/2}\,\E_{y\sim N(0,\tau I)}\,\delta(y)\,d\tau ,
\end{equation}
and the weight $\tfrac12(1+\tau)^{-3/2}d\tau$ is a probability measure on $[0,\infty)$. At the origin $\delta(0)=e(0)-\Delta
e(0)$, and the exact answer satisfies $u=\Delta_y u(y)|_{y=0}$ (Euler--Stein).
\end{corollary}

\begin{proof}
On Eldan's path $\nu_t=N(m_t,s_tI)$, $s_t=1/(1+t)$, $ds_t=-s_t^2dt$, $d\langle m\rangle_t=s_t^2I\,dt$. With
$E(m,sI)=\sqrt s\,e(m/\sqrt s)$ one finds $\partial_sE=\tfrac1{2\sqrt s}(e-y\cdot\nabla e)$ and $\tfrac12\Delta_mE=
\tfrac1{2\sqrt s}\Delta e$ at $y=m/\sqrt s$, so the drift is $-s^2(\partial_sE-\tfrac12\Delta_mE)=-\tfrac12 s^{3/2}\delta(y)$.
Since $y_t=m_t\sqrt{1+t}\sim N(0,tI)$, Theorem~\ref{thm:defect} gives \eqref{eq:eldan} with $\tau=t$. For the exact
answer, $u(y)=\E F(y+Z)$ obeys the heat equation $\partial_s u=\tfrac12\Delta_mu$, and for a $1$-homogeneous $F$ Euler's
relation $x\cdot\nabla F=F$ and Stein's lemma $\E[x\cdot\nabla F(x)]=\E\,\Delta F(x)$ give $u=\Delta_yu|_0$.
\end{proof}

This corollary is the heat-flow defect of the parallel note (its ``heat defect'' is $D=\partial_\Sigma E-\tfrac12\nabla_m^2E$;
here $\tfrac12 s^{3/2}\delta$ is the trace of $D\Sigma_t^2$ on the isotropic path). Theorem~\ref{thm:defect} is the
scheme-independent statement behind it, and the choice of scheme is where the stage-8 picture adds something.

\begin{corollary}[Directional and pinning schemes]\label{cor:directional}
(i) For the directional tilt along a unit vector $v$, the path is $\nu_t=N(m_t,I-\tfrac{t}{1+t}vv^\top)$ with
$d\langle m\rangle_t=(1+t)^{-2}vv^\top dt$, the endpoint is $\nu_\infty=N(\langle v,X\rangle v,\,I-vv^\top)$, and
\[
\E_{y}\,E\big(N(yv,I-vv^\top)\big)-E(N(0,I))\;=\;-\int_0^\infty(1+t)^{-2}\,\E\Big[\tfrac12\partial_v^2E-\partial_{\Sigma_{vv}}E\Big](\nu_t)\,dt ,
\]
the left side being the error of the $v$-localized estimator minus the error of $E$. So the drift along $v$ measures
the \emph{difference} between the chain and its $v$-localization, and the chain's own error is recovered from
directional drifts only by summing over an orthonormal basis of directions---which is Eldan's path.
(ii) For pinning (reveal $X_i$ in order), the Jensen gap at step $i$ is
$\E_{X_i}E(\nu_{i-1}\mid X_i)-E(\nu_{i-1})$, a one-dimensional Gaussian integral of the chain run on the conditioned
law; the chain's error is minus the sum of these $n$ gaps.
\end{corollary}

\begin{remark}[Which errors a scheme sees]\label{rem:visibility}
The second variation $E''(\nu)[\dot\nu,\dot\nu]$ is a quadratic form on the tangent space of $\mathcal S_0$ at $\nu$.
A scheme supplies a positive form $d\langle\nu\rangle$ (its driving covariance $C_t$ in the language of
\cite[Thm.~3.3]{stage8}, which also fixes the trickle-down term $KC_tK$ of that theorem). Eldan's path drives every
mean direction equally and the covariance only through the scale; the directional tilt drives one direction; pinning
drives the coordinate directions and their conditional variances. Which scheme probes a given chain cheaply is
therefore a question about the \emph{spectrum} of the chain's second variation on the input side. Section~\ref{sec:capacity}
measures it: for the chains of this note the excess second variation $\partial_v^2e-\Delta e/n$ over unit input
directions $v$ has a flat spectrum (the top direction carries $5\%$ of it at $n=256$, and is not the coherent
direction $W_1^{\top}\mathbf 1$), and single directional second differences are nearly uncorrelated with the error.
The chain's error is an isotropic, trace-like object on the input side---the Laplacian of Corollary~\ref{cor:eldan}---and
the only cheap way to get at it is not a tilt of the input but a second-order expansion \emph{inside} the chain
(Section~\ref{sec:central}). The collective mode that the chain mishandles lives in the gate algebra of the hidden
layers, not in the input, which is the content of the capacity bound of Section~\ref{sec:capacity}.
\end{remark}

\section{Births are linear-tilt responses of the gate algebra}
\label{sec:births}
We now compute the second variation of a chain explicitly. The result is that the objects a cumulant chain
transports (births at the kinks, gated arms, $(2,1)$ slices) are the first and second derivatives of the Gaussian
moment map under a shift of the mean, and that the Euler--Stein defect of Section~\ref{sec:capacity} is the
curvature of the downstream chain contracted with the Gram matrix of those derivatives.

\subsection{The chaos index shifts under a tilt}
For $u\sim N(0,1)$ and $a\in\R$ let $r_a(u)=\relu(u+a)$ with Hermite coefficients $d_k(a)=\E[r_a(u)\Hpl_k(u)]$:
\[
d_0=\varphi(a)+a\Phi(a),\qquad d_1=\Phi(a),\qquad d_k=\varphi(a)\,\Hpl_{k-2}(-a)\ (k\ge2),
\]
so that for standardised $(u,v)$ with correlation $\rho$, $\E[r_a(u)r_b(v)]=\sum_{k\ge0}\rho^kd_k(a)d_k(b)/k!$.
\begin{lemma}[Index shift]\label{lem:shift}
$d_k'(a)=d_{k+1}(a)$ for every $k\ge0$.
\end{lemma}
\begin{proof}
$d_0'=-a\varphi+\Phi+a\varphi=\Phi=d_1$ and $d_1'=\varphi=d_2$. For $k\ge2$, with $x=-a$,
$\frac{d}{da}[\varphi(a)\Hpl_{k-2}(x)]=\varphi\,[x\Hpl_{k-2}(x)-(k-2)\Hpl_{k-3}(x)]=\varphi\,\Hpl_{k-1}(x)=d_{k+1}(a)$
by $\Hpl_m'=m\Hpl_{m-1}$ and the three-term recursion $\Hpl_{m+1}(x)=x\Hpl_m(x)-m\Hpl_{m-1}(x)$.
\end{proof}
(Checked numerically to $2\cdot10^{-8}$ for $k\le10$ in \texttt{verify\_stage9.py}.) The lemma is Stein's identity
$\E[f'(u)\Hpl_k(u)]=\E[f(u)\Hpl_{k+1}(u)]$ for $f=r_a$; what matters here is that the tilt $a\mapsto a+t$ of the
Gaussian acts on the chaos expansion of the \textsc{ReLU} by raising the index.

Let now $z\sim N(m,S)$ be the pre-activation of one layer, $\sigma_a=\sqrt{S_{aa}}$, $\alpha_a=m_a/\sigma_a$,
$\rho_{ab}=S_{ab}/\sigma_a\sigma_b$, and
\[
\mu_a=\sigma_ad_0(\alpha_a),\qquad M_{ab}=\E[\relu z_a\relu z_b]=\sigma_a\sigma_b\sum_k\rho_{ab}^k\,d_k(\alpha_a)d_k(\alpha_b)/k!
\]
the post-activation mean and raw second moment. Lemma~\ref{lem:shift} gives every derivative of the closure:
\begin{align}
\partial_{m_a}\mu_a&=\Phi(\alpha_a)=:g_a, & \partial_{m_a}^2\mu_a&=\varphi(\alpha_a)/\sigma_a=:w_a,\label{eq:birth}\\
\partial_{m_a}M_{ab}&=\sigma_b\sum_k\rho^kd_{k+1}(\alpha_a)d_k(\alpha_b)/k!=\E[\theta(z_a)\relu z_b]=:T_{ab},\label{eq:slice}\\
\partial_{S_{ab}}M_{ab}&=\sum_k\rho^kd_{k+1}(\alpha_a)d_{k+1}(\alpha_b)/k!=\E[\theta(z_a)\theta(z_b)]=\partial_{m_a}\partial_{m_b}M_{ab}=:G_{ab},\label{eq:gate}\\
\partial_{m_a}^2M_{ab}&=\tfrac{\sigma_b}{\sigma_a}\sum_k\rho^kd_{k+2}(\alpha_a)d_k(\alpha_b)/k!=w_a\mu_b+\tfrac{\sigma_b}{\sigma_a}\sum_{k\ge1}\rho^kd_{k+2}(\alpha_a)d_k(\alpha_b)/k!.\label{eq:birth2}
\end{align}
Here $\theta$ is the Heaviside gate. The identities with expectations are Price's theorem
$\partial_{S_{ab}}\E[f(z_a)h(z_b)]=\E[f'(z_a)h'(z_b)]$ and Stein's lemma; the closure is exact for one layer
because it satisfies the heat equation $\partial_{S_{ab}}=\partial_{m_a}\partial_{m_b}$ in its own parameters.
The dictionary with the cumulant chain of \cite{Wu} and of Section~\ref{sec:transport} is now literal:
$w_a=\varphi(\alpha_a)/\sigma_a$ is the \emph{birth} strength at the kink of neuron $a$, the second derivative of
the mean under a tilt; $g_a=\Phi(\alpha_a)$ gates the arms; $T_{ab}$ is the exact $(2,1)$ slice function
$\E[\theta(z_a)\relu z_b]$ through which a tilt at $a$ moves the covariance row of $a$; $G_{ab}$ is the gate
covariance $\E[\theta_a\theta_b]$ that transports a covariance perturbation. Births are the curvature of the
tilt response; arms and slices are its gradient.

\subsection{The defect is downstream curvature against the tilt Gram}
Write the chain as $e(y)=\mathcal G(\theta(y))$, where $\theta(y)=(\mu(y),M(y))$ are the exact first-layer
post-activation moments for the input law $N(y,I)$ (exact because $z_1=W_1x$ is Gaussian) and $\mathcal G$ is the
chain from layer $2$ on, a function of the Gaussian surrogate $N(\mu,M-\mu\mu^\top)$ of $h_1$.
\begin{theorem}[Defect $=$ curvature $\times$ Gram]\label{thm:gram}
Let $\mathcal G$ be homogeneous, $\mathcal G(\lambda\mu,\lambda^2M)=\lambda\,\mathcal G(\mu,M)$, and twice
differentiable. Then for every $y$,
\begin{equation}\label{eq:gram}
\delta(y)\;=\;e-y\cdot\nabla e-\Delta e\;=\;-\sum_{A,B}\partial_A\partial_B\mathcal G\big(\theta(y)\big)\,\Gamma_{AB}(y),
\qquad \Gamma_{AB}(y)=\big\langle\nabla_y\theta_A(y),\nabla_y\theta_B(y)\big\rangle ,
\end{equation}
the sum running over the moment indices $A,B\in\{\mu_a\}\cup\{M_{ab}\}$. With $S=W_1W_1^\top$ and the rows $w_a$ of
$W_1$, at $y=0$ the Gram matrix is
\[
\Gamma_{\mu_a\mu_b}=g_ag_bS_{ab},\qquad
\Gamma_{\mu_aM_{bc}}=g_a\big(T_{bc}S_{ab}+T_{cb}S_{ac}\big),\qquad
\Gamma_{M_{ab}M_{cd}}=T_{ab}T_{cd}S_{ac}+T_{ab}T_{dc}S_{ad}+T_{ba}T_{cd}S_{bc}+T_{ba}T_{dc}S_{bd}.
\]
Consequently $u-e(0)=\int_0^\infty\!w(\tau)\,\E_{y\sim N(0,\tau I)}\langle\nabla^2\mathcal G(\theta(y)),\Gamma(y)\rangle\,d\tau$
with $w(\tau)=\tfrac12(1+\tau)^{-3/2}$.
\end{theorem}
\begin{proof}
Each $\theta_A(y)=\E\,g_A(y+Z)$, $Z\sim N(0,I)$, with $g_A$ positively homogeneous of degree $p_A\in\{1,2\}$,
satisfies the heat equation $\partial_s\theta_A=\tfrac12\Delta\theta_A$ along $N(y,sI)$ and the scaling
$\theta_A(y,s)=s^{p_A/2}\theta_A(y/\sqrt s,1)$, whence $p_A\theta_A-y\cdot\nabla\theta_A-\Delta\theta_A=0$. By the
chain rule $\nabla e=\sum_A\partial_A\mathcal G\,\nabla\theta_A$ and $\Delta e=\sum_A\partial_A\mathcal G\,\Delta\theta_A
+\sum_{A,B}\partial_A\partial_B\mathcal G\,\Gamma_{AB}$, so
$\delta=\mathcal G-\sum_Ap_A\theta_A\partial_A\mathcal G-\sum_{AB}\partial_A\partial_B\mathcal G\,\Gamma_{AB}$, and
Euler's relation for the weighted homogeneity of $\mathcal G$ kills the first two terms. The Gram entries follow
from \eqref{eq:birth}--\eqref{eq:slice}: $\nabla_y\mu_a=g_aw_a$ and $\nabla_yM_{ab}=T_{ab}w_a+T_{ba}w_b$. The last
statement is Corollary~\ref{cor:eldan}.
\end{proof}

Three remarks. (1) The theorem is exact and holds for every homogeneous chain (Gaussian closure or any cumulant
chain run from the Gaussian surrogate of $h_1$); nothing about the chain's internals enters except its curvature
in the first-layer moments. (2) The Gram matrix has the \emph{CP-leg} form of the cumulant chain: a gated
covariance $g\circ S\circ g$ (the arms), slice products $T\circ S$ (the $(2,1)$ legs), and the four-term
$(M,M)$ block in which the pair $(a,b)$ is linked to $(c,d)$ through a single covariance entry. The same
construction at a hidden layer $l$ (tilt the hidden mean $m_l$ with the chain's $C_l$ in place of $I$) gives the
error of the chain started from the surrogate at layer $l$, $\mathrm{Err}_l$, and the differences
$\mathrm{Err}_{l-1}-\mathrm{Err}_l$ are the injected errors of Section~\ref{sec:measure}. (3) The contraction
$\langle\nabla^2\mathcal G,\Gamma\rangle$ can be evaluated inside the chain by propagating the second-order
response of the closure to a Gaussian shift of the mean, i.e.\ the covariances $\E[\delta m\,\delta m^\top]$,
$\E[\delta m_a\,\delta S_{ab}]$, $\E[\delta m_a\,\delta S_{bb}]$, $\E[\delta S_{aa}\delta S_{bb}]$,
$\E[\delta S_{ab}\delta S_{aa}]$ of the perturbation; every one of them closes at $O(n^3)$ per layer through
\eqref{eq:birth}--\eqref{eq:birth2} except $\E[\delta S_{ab}^2]$, the variance of the transported
cross-covariance, whose exact propagation is $O(n^4)$. This is the same obstruction that forces a cumulant chain
to carry slices of $\kap_3$ rather than the tensor, met here from the other side: the second variation of a
chain has the complexity of the chain's own closure error.

\section{Capacity of input tilts}
\label{sec:capacity}
The stage-8 programme asks first which localization sees the error. Here we settle it for the input side, with
two measurements on a small network ($n=256$, $L=4$, Monte Carlo truth from $10^7$ samples, noise $2\cdot10^{-4}$)
and one on the official network $0$ ($n=1024$, $L=16$).

\subsection{The Laplacian is the error; its directional parts are not}
Table~\ref{tab:defect-small} gives, for the Gaussian closure and for the chain of this note (\texttt{k3v3}), the
per-neuron correlation between the error $e(0)-u$ and the Euler--Stein defect $\delta(0)=e(0)-\Delta e(0)$ of
Corollary~\ref{cor:eldan}, the least-squares coefficient $a$ in $e(0)-u\approx a\,\delta(0)$, and the mean squared
error after the correction $e(0)-a\delta(0)$. The Laplacian was computed exactly by $2n+1$ chain evaluations.
\begin{table}[h]\centering\small
\begin{tabular}{lcccccc}\toprule
chain & MSE & $|\delta|_{\rm rms}$ & corr$(e-u,\delta)$ & $a$ & fraction explained & MSE after $e-a\delta$\\\midrule
Gaussian closure & $1.16\cdot10^{-5}$ & $9.7\cdot10^{-3}$ & $0.967$ & $0.339$ & $0.936$ & $7.5\cdot10^{-7}$\\
\texttt{k3v3}    & $1.09\cdot10^{-6}$ & $6.1\cdot10^{-3}$ & $0.835$ & $0.143$ & $0.690$ & $3.4\cdot10^{-7}$\\
\bottomrule\end{tabular}
\caption{Euler--Stein defect against the true error, $n=256$, $L=4$ (\texttt{scripts/defect\_small.py}).}
\label{tab:defect-small}
\end{table}
Two facts. First, $a<1$: Corollary~\ref{cor:eldan} integrates $\delta(y)$ against $y\sim N(0,\tau I)$ with the
weight $\tfrac12(1+\tau)^{-3/2}$, and $\delta(y)$ decays as $|y|$ grows (a shifted input makes the network more
nearly deterministic relative to its fluctuations), so $\delta(0)$ overstates the integral; $a$ is the ratio of
the weighted mean of $\delta$ to its value at the origin, a property of the architecture class rather than of the
network, and is fitted once. Second, the \texttt{k3v3} chain, which carries the third cumulants, has a smaller
$a$ and a smaller explained fraction: its defect is of higher order and less self-similar in $y$.

The directional second differences tell the opposite story. For a unit vector $v$ let $\partial_v^2e$ be the
second difference of $e$ along $v$ at $0$ and $x_v=\partial_v^2e-\Delta e/n$ its excess over the isotropic share.
On the same network the singular values of the matrix $(x_{e_i})_{i\le n}$ decay slowly (top share $5\%$, top four
$17\%$), the leading direction is orthogonal to the coherent direction $W_1^{\top}\mathbf 1$ ($|\langle u_1,
v_{\rm coh}\rangle|<10^{-3}$), and the correlation between $x_v$ and the error is $0.18$, $0.12$, $0.08$
(coherent, random, leading) for the Gaussian closure and $0.26$, $-0.05$, $-0.15$ for \texttt{k3v3}. No input
direction carries the error.

\subsection{The defect is a small residual of an $O(1)$ identity}
At $n=1024$, $L=16$ (official network $0$) the random-probe estimate of the Laplacian,
$\Delta e\approx\frac1S\sum_s\partial_{\xi_s}^2e$ with $\xi_s\sim N(0,I)$ (Hutchinson), fails outright: with
$S=16$ probes per block the fitted coefficient is $a\approx0.01$ and the correlation with the error changes sign
from block to block ($+0.43,-0.09,-0.44,+0.14$). The reason is quantitative. Euler--Stein says $\Delta u=u$ for
the exact answer, and the chain nearly obeys it: $\Delta e(0)=e(0)-\delta(0)$ with $|\delta|\approx10^{-2}$ against
$|e|\approx0.4$. The quantity we need, $\delta$, is a $2\%$ residual of an $O(1)$ identity, and the error is a
tenth of that residual, so the Laplacian must be computed to a relative precision of $10^{-4}$. For a unit
direction $v$ one measures $|\partial_v^2e|_{\rm rms}\approx4\cdot10^{-3}$ at $n=256$; since
$\E_v(\partial_v^2e)^2-(\tr H/n)^2=2\|H\|_F^2/n^2$ for $v$ uniform on the sphere, the Hessian $H=\nabla^2e$ of one
output neuron has $\|H\|_F\approx n|\partial_v^2e|/\sqrt2\approx0.8$, comparable with $\tr H\approx0.4$. The
variance of one Hutchinson probe is $2\|H\|_F^2$, so the relative noise of an $S$-probe estimate of $\tr H$ is
$\sqrt{2/S}\,\|H\|_F/\tr H\approx2.8/\sqrt S$, and the required $10^{-4}$ needs $S\approx10^9$ probes. Only the
exact coordinate sum, $2n+1$ chain evaluations, resolves $\delta$.

\begin{proposition}[Capacity of input tilts]\label{prop:capacity}
Let $E$ be a chain exact on point masses and homogeneous as in Corollary~\ref{cor:eldan}, and let $H$ be the
Hessian of $e$ at $0$ for one output neuron. (i) The error is the trace part: $u-e(0)=-\int w(\tau)\E_y\delta(y)\,d\tau$
with $\delta(0)=e(0)-\tr H$. (ii) A localization that drives $k$ orthonormal directions $v_1,\dots,v_k$ (directional
tilts, or pinning of $k$ coordinates) has drift $-\sum_{s\le k}(\tfrac12\partial_{v_s}^2E-\partial_{\Sigma_{v_sv_s}}E)$.
Over a random choice of the $k$ directions its expectation is $(k/n)$ times the isotropic drift, whose
error-carrying part is $(k/n)\,\delta$, while its fluctuation over the choice is of order $\sqrt{k}\,\|A\|_F/n$
with $A=\tfrac12H-\partial_\Sigma E$. The signal-to-noise ratio of a $k$-direction scheme is therefore
$\sqrt k\,|\delta|/\|A\|_F$.
(iii) With $|\delta|/\|A\|_F\approx10^{-2}$ as measured, a $k$-direction scheme resolves the error only for
$k\gtrsim10^4\gg n$, and the cheapest exact scheme on the input side is the coordinate sum, $2n+1$ chain
evaluations, i.e.\ $2n+1$ pair chains at $n=1024$, far outside any budget. Input localization is useless for
estimation; it is a diagnostic.
\end{proposition}

\begin{proof}
(i) is Corollary~\ref{cor:eldan}. (ii): the directional drift of Corollary~\ref{cor:directional} is a quadratic
form in $v$ with matrix $\tfrac12H-\partial_\Sigma E$, whose average over the Haar measure on $k$-frames is $k/n$ of
its trace, and whose variance over the frame is $\sum_s\Var(v_s^{\top}Av_s)\approx 2k\|A\|_F^2/n^2$ for $A$ the
matrix of the form. (iii) follows by comparing the two.
\end{proof}

This is the negative half of the stage-8 dictionary: the input is the \emph{commutative} side, its tilts are the
modular tilts of the Gaussian, and the chain's defect is invisible to any finite family of them because it is a
trace. The positive half is that the trace can be computed \emph{inside} the chain, as the second-order response
of the closure to a Gaussian shift of the mean, which is the subject of Sections~\ref{sec:births} and
\ref{sec:central}: there the directions are neurons, the tilt responses are the births at the kinks, and the sum
over directions is a sum over neurons that the chain already performs.

\section{The collective mode is the central variable}
\label{sec:central}
The stage-8 principle splits a localization into its trivial (central) and Cartan parts. For the \textsc{MLP} the
central variable is the collective activity of a layer: the direction $s=\E[z^2]$ in pre-activation space, which is
also the top eigenvector of the covariance once depth has built it up. This section proves two small exact
statements about it and reports what the chain does to it.

\subsection{A scale mixture is repaired exactly at first order by the $(2,2)$ channel}
\begin{lemma}[Scale mixture]\label{lem:scale}
Let $z\mid\varepsilon\sim N(m,(1+\varepsilon)^2S)$ with $\E\varepsilon=0$, $\E\varepsilon^2=v$ small. Then the
cumulants of $z$ are those of $N(m,(1+v)S)$ at second order and $\kap_4(z_a,z_b,z_c,z_d)=4v\,(S_{ab}S_{cd}+
S_{ac}S_{bd}+S_{ad}S_{bc})+O(v^2)$, a pure $(2,2)$-type tensor. The Gaussian closure applied to the true
covariance $(1+v)S$ gives $\E\relu z_a$ with error $-\tfrac v2\sigma_a(\alpha_a^2-1)\varphi(\alpha_a)+O(v^2)$, and
the $\kap_4$ term of the mean readout \eqref{eq:mean-readout} with $\kap_4(z_a)=12v\sigma_a^4$ is
$+\tfrac v2\sigma_a(\alpha_a^2-1)\varphi(\alpha_a)$: the two cancel.
\end{lemma}
\begin{proof}
$\E[(1+\varepsilon)^4]=1+6v+O(v^2)$ and $\E[(1+\varepsilon)^2]^2=1+2v+v^2$ give the fourth cumulant. Put
$f(\varepsilon)=(1+\varepsilon)\sigma d_0(\alpha/(1+\varepsilon))$; the true mean is $\E f(\varepsilon)=f(0)+\tfrac v2f''(0)$
and the closure's is $f(\varepsilon^*)$ with $(1+\varepsilon^*)^2=1+v$, i.e.\ $f(0)+\tfrac v2f'(0)$. From
Lemma~\ref{lem:shift}, $f'(0)=\sigma[d_0-\alpha d_1]=\sigma\varphi$ and $f''(0)=\sigma\alpha^2\varphi$.
\end{proof}
So the $(2,2)$ channel of the chain (\texttt{k22}: $\kap_4(h_a,h_a,h_b,h_b)$ born exactly from the Hermite series of
$\Cov((h_a-\mu_a)^2,(h_b-\mu_b)^2)$, transported by $(W\!\circ\!W)\,K\,(W\!\circ\!W)^\top$ with the gate $tt^\top$ and
the diagonal class removed once) is the first-order correction for a fluctuating collective scale. Its gain
projection $4v=s^\top K_{22}s/(s^\top s)^2$ measures $v$ directly.

\subsection{The sensitivity law}
\begin{proposition}[Sensitivity]\label{prop:sens}
A relative error $\varepsilon_a$ in the pre-activation variance of neuron $a$ changes its post-activation mean by
$\tfrac12\varepsilon_a\sigma_a\varphi(\alpha_a)$, and an error $\delta\lambda$ in the top eigenvalue of the
covariance of layer $l$ (eigenvector $u$) changes the pre-activation variance of every neuron $b$ of layer $l+1$
by $\delta\lambda\,(w_b\!\cdot\!u)^2\approx2\,\delta\lambda/n$. On the official networks ($\sigma\approx2.8$,
$\varphi(\alpha)\approx0.3$, $\mathrm{var}\approx8$, $\lambda_1\approx12$ at depth) a uniform relative variance
error of $10^{-3}$ moves every mean by $4\cdot10^{-4}$, seven times the rms error of the chain, while a $2\%$ error of
the top eigenvalue moves the next layer's variances by $5\cdot10^{-5}$ relative and the means by $2\cdot10^{-5}$ per
layer.
\end{proposition}
\begin{proof}
$\partial_{\sigma^2}[\sigma d_0(m/\sigma)]=\tfrac1{2\sigma}(d_0-\alpha d_1)=\varphi/(2\sigma)$, multiplied by
$\varepsilon\sigma^2$; and $w_b^\top(\delta\lambda\,uu^\top)w_b$ with $w_b\!\cdot\!u\sim N(0,2/n)$ for He rows.
\end{proof}
The bulk of the covariance may therefore be wrong by percents without consequence (it enters the next layer's
variances through $\sum_cS_{bc}^2$-type sums that average out), while the diagonal must be right to $10^{-4}$ and
the collective mode to $10^{-3}$: that is what makes the central variable the one that matters.

\subsection{Measurements on the official network 0}
From the Monte Carlo cache ($2^{20}$ inputs; \texttt{scripts/diag\_eig\_v3.py}, \texttt{diag\_onestep.py}):
\begin{itemize}[nosep]
\item The collective mode emerges with depth: the top four eigenvalues of the pre-activation covariance are
$8.08,7.89,7.77,7.71$ at layer $0$ and $16.3,5.36,4.88,4.62$ at layer $15$.
\item The chain's trace is exact to $10^{-3}$ at every layer and its diagonal to $10^{-4}$; the bulk (the
complement of the top four modes, off-diagonal part) is off by $2$--$4\%$; the top eigenvalue is correct to
$2\cdot10^{-3}$ through layer $5$ and then $1.0,1.2,1.5,1.8,1.9,1.8,1.7,1.8,1.5,1.6\%$ \emph{low} at layers
$6$--$15$, with the variance along $s$ equally low. By Proposition~\ref{prop:sens} this is a $2\cdot10^{-5}$ per
layer effect on the means, a tenth of the measured error: the top mode is the \emph{symptom} of the deficit, not
its carrier.
\item One-step closure from Monte Carlo inputs at layers $11,13,14$: the Gaussian closure loses $3.2,1.5,1.1\%$ of the
top mode in one layer; adding the Monte Carlo $(2,1)$ slices (the $\kap_3$ covariance correction
\eqref{eq:cov-transport}) leaves $-0.13,-0.09,-0.09\%$; adding the $(2,2)$ slices changes that by less than
$3\cdot10^{-4}$; adding the $(3,1)$ slices makes it worse ($-0.19,-0.15,-0.16\%$). The $\kap_3$ slice is the whole
one-step story of the collective mode, and the sign of the $(3,1)$ term says its first-order transport is not the
right object to add.
\item The chain's own transported slices at those layers have correlation $0.965$ with Monte Carlo and slope
$0.92,0.90,0.91$; its $(2,2)$ slices have correlation $0.22$--$0.25$ and slope $0.53$--$0.68$, and its gain
projection $4v$ is $3.6,1.9,1.5\cdot10^{-5}$ against $6.6,4.6,4.0\cdot10^{-5}$: the scale-mixture variance is $v\approx
10^{-5}$, so by Lemma~\ref{lem:scale} the $(2,2)$ channel is worth $\sim10^{-5}$ per neuron and layer, which is
why the chain's $(2,2)$ errors are harmless.
\end{itemize}
The error accounting that follows is then simple. A $10\%$ deficit of the transported slice leaves $10\%$ of the
$1$--$3\%$ one-step loss of the collective mode in place, $0.1$--$0.3\%$ per layer, which accumulates to the
$1.5$--$1.9\%$ measured; the same deficit leaves $10\%$ of the $\kap_3$ corrections to the diagonal and the means in
place at every layer, and the per-layer injected mean errors of $3\cdot10^{-4}$ (Section~\ref{sec:measure}) are of
that size: $10\%$ of $\kap_3(z_a)\alpha\varphi/(6\sigma^2)$ at $\kap_3\approx3$, which is $3\cdot10^{-4}$. The chain's
error is the deficit of its $(2,1)$ transport, in the diagonal readout and in the collective mode alike.

\section{Exact first-order transport through \textsc{ReLU}}
\label{sec:transport}
The chain of this note (\texttt{whest/k3chain3.py}) transports the third cumulant in the CP-leg form of
\cite{Wu}: each kink neuron $m$ of layer $l$ is a source with centre leg $P=W\!\cdot e_m$, arm $A=W\!\cdot a_m$,
$a_m=\diag(g)\,Ce_m$, strength $w_m=\varphi(\alpha_m)/\sigma_m$, and the two further legs $B$ (the Gram--Charlier
arm fed by the incoming $(2,1)$ slices) and $R$ (the residual that makes the slice and the diagonal exact at
birth). The readouts in pre-activation space are the diagonal $\kap_3(z_a)=D3_a$ and the slices
$D21_{ab}=\kap_3(z_a,z_a,z_b)$. What is new relative to \cite{Wu} is the \emph{transport} of a slice through a
\textsc{ReLU} layer to first order in the cumulants, in closed form, which we state and then measure.

\begin{theorem}[First-order transport of the $(2,1)$ slice]\label{thm:transport}
Let $z\sim$ a law with Gaussian part $N(m,S)$ and small third and fourth cumulants, $h=\relu(z)$, and write
$g=\Phi(\alpha)$, $w=\varphi(\alpha)/\sigma$, $\mu_a=\sigma_ad_0(\alpha_a)$, $t_a=g_a-\mu_aw_a$,
$\epsilon_a=2\mu_a(1-g_a)$, $c_a=g_a(1-g_a)-\mu_aw_a$. To first order in the cumulants of $z$ and to the orders in
$\rho$ kept by the chain,
\begin{align}
\kap_3(h_a,h_a,h_c)&=\Big[g_a^2+c_a\Big]\,g_c\,\kap_3(z_a,z_a,z_c)
+\tfrac12\,\epsilon_a\,w_c\,\kap_3(z_a,z_c,z_c)+\tfrac12\,t_a\,w_c\,\kap_4(z_a,z_a,z_c,z_c)+\text{(born)},\label{eq:slice-transport}\\
\Cov(h_a,h_b)&=\Cov^{\rm Gauss}(h_a,h_b)+\tfrac12\big(w_ag_b\,\kap_3(z_a,z_a,z_b)+g_aw_b\,\kap_3(z_a,z_b,z_b)\big)
+\tfrac14w_aw_b\,\kap_4(z_a,z_a,z_b,z_b),\label{eq:cov-transport}\\
\E[h_a]&=\sigma_ad_0(\alpha_a)-\tfrac{\alpha_a\varphi(\alpha_a)}{6\sigma_a^2}\kap_3(z_a)+\tfrac{(\alpha_a^2-1)\varphi(\alpha_a)}{24\sigma_a^3}\kap_4(z_a),\label{eq:mean-readout}
\end{align}
where ``born'' is the exact $\kap_3$ of $\relu$ of the Gaussian part (all Hermite orders of the slice,
\texttt{relu2\_hermite}, and the exact diagonal \texttt{k3\_exact}), and the gate of the transported slice,
$g_a^2+c_a=g_a-\mu_aw_a$, is the exact coefficient $\E[\theta_a\,\partial_a(\relu z_a-\mu_a)^2]/2$ rather than
the Gaussian $g_a^2$.
\end{theorem}
\begin{proof}[Proof sketch]
Expand the law of $z$ as a Gram--Charlier perturbation of its Gaussian part, $\nu=\nu_G(1+\tfrac16\kap_3\!:\!\Hpl_3
+\tfrac1{24}\kap_4\!:\!\Hpl_4)$, and compute $\E[(h_a-\mu_a)^2(h_c-\mu_c)]$ by integrating the Hermite
polynomials against the \textsc{ReLU} moments: Stein's lemma turns each $\Hpl_3$ against a product of
\textsc{ReLU}s into third derivatives, which are sums of gate products $\theta_a\theta_a\theta_c$ (the first term),
of $\theta_a\delta_c\relu_c$-type terms, i.e.\ a birth $w_c$ at $c$ against the transposed slice (the second term),
and the $\Hpl_4$ term gives the $(2,2)$ feed with the single birth $w_c$ (third term). The coefficients
$t_a,\epsilon_a,c_a$ are the exact one-dimensional integrals $\E[\theta_a(\relu z_a-\mu_a)\cdot]$ listed in the code.
The covariance and mean formulas are the same computation one order lower.
\end{proof}

The three transported terms are what the chain calls the gated slice (\texttt{gate\_res}), the transposed-slice
feed (\texttt{feedT}) and the $\kap_4\!\to\!\kap_3$ feed (\texttt{feedK22}). Their effect on the official
network $0$ is the ladder of Section~\ref{sec:measure}: the transported slices rise from $75\%$ to $90$--$92\%$ of
their Monte Carlo values at depth and the final MSE falls from $8.5\cdot10^{-7}$ to $5.7\cdot10^{-7}$.

\paragraph{What remains open in the transport.} The first term of \eqref{eq:slice-transport} transports the slice
with the gate of the pair $(a,c)$ taken at $\rho_{ac}=0$ plus the chain's $\rho$-expansion; the exact gate
$\E[\theta_a\theta_a\theta_c\,(\cdot)]$ at the true correlation differs at order $\rho_{ac}$, which is small per entry
but coherent along the collective mode. Section~\ref{sec:dictionary} lists the $\rho$-exact gate as the first
open item, since the $8$--$10\%$ deficit of the transported slice at depth is precisely of this size.

\section{What the chain costs: the pair floor}
\label{sec:cost}
The score multiplies the error by $\max(0.1,\,C/B)$ with $C$ the metered count of the grader's own numpy wrapper
(\texttt{flopscope} 0.12.1) and $B=2^{41}$. We read the wrapper's rules off its source and measured the bills of
the chain's building blocks at $n=1024$ (\texttt{notes/compute/FLOPSCOPE.md}); the facts that decide a design are
these.

\paragraph{Rules.} Every call bills $\lfloor\text{textbook count}\times\text{dtype rate}\times\text{weight}\rfloor$.
The count is the textbook one, not what BLAS does: a matmul $(m,k)(k,n)$ is $2mkn-mn$ whatever algorithm the
library uses, so a Strassen level written out of seven smaller \texttt{matmul} calls \emph{is} billed at $7/8$.
The dtype rate is $1$ for float32 and $2$ for float64. The weight is $0$ for views, transposes, \texttt{zeros};
$1$ for arithmetic, reductions, contractions, copies, \texttt{sqrt}, \texttt{stats.norm.cdf/pdf} (whose counts are
$48$ and $27$ per element, computed in float64), \texttt{polyval} ($2d$ per element) and all of \texttt{linalg};
$4$ for \texttt{where}, gathers and sorts; $16$ for \texttt{exp}, \texttt{log} and every random sampler. A
contraction whose output the engine can prove symmetric is billed on the unique entries: the sandwich
$WCW^\top$ written as one \texttt{einsum('ij,jk,lk->il', W, C\_sym, W)} with the same array $W$ in both outer
slots and $C$ tagged \texttt{as\_symmetric} bills $3n^3$ in float32 ($6n^3$ in float64) against $4n^3$ ($8n^3$) for
two matmuls; the saving is lost if the two products are separate calls. \texttt{eigh} is $9n^3$, \texttt{cholesky}
$n^3/3$, \texttt{inv} $2n^3$, \texttt{solve} $\tfrac23n^3+2n^2r$.

\paragraph{Budget.} The floor $0.1B=2.2\cdot10^{11}$ equals $205\,n^3$ over the $16$ layers, i.e.
$12.8\,n^3$ per layer: six float32 $(n,n,n)$ matmuls, or four float32 symmetric sandwiches, or three float64
matmuls. Nothing is gained below the floor, so a chain should spend exactly it. The Gaussian pair chain,
one sandwich per layer, costs $0.024B$ in float32 (one quarter of the floor) and $0.047B$ in float64.
Elementwise work on $n\times n$ arrays is negligible unless repeated dozens of times per layer: the Hermite series
of the pair moment to order $K$ is $K$ Hadamard products, $2Kn^2$, i.e.\ $0.002\,n^3$ per term.

\begin{proposition}[Pair floor]\label{prop:floor}
Any estimator that forms the pre-activation covariance of every layer from that of the previous one pays at least
one sandwich per layer: $3n^3$ in float32 with the symmetric engine, $L\cdot3n^3=0.024B$ at $n=1024$, $L=16$.
Every further $n\times n$ object transported through $W$ (a $\kap_3$ slice, a $(2,2)$ slice, a gate-correlation
matrix) costs another $2n^3$ to $3n^3$ per layer, i.e.\ $0.016B$ to $0.024B$ over the network. Under the floor a
chain may therefore carry the pair plus about four such objects per layer; the leaders' bills ($0.105B$,
$0.144B$) correspond to five to seven float32 sandwich-equivalents per layer, and the open-source third-cumulant
chain at $0.25B$ to a float64 implementation of the same count.
\end{proposition}

\paragraph{Naive counts mislead.} The count our chain returns (FMA $=2$, no dtype rate, no symmetry) overstates
the float32 symmetric bill by $2.7\times$ and understates a float64 plain-matmul bill by $2\times$. The
\emph{pair floor} and the \emph{per-object price} above are the only two numbers a design needs; every proposal in
Section~\ref{sec:dictionary} is costed in units of ``$n\times n$ objects transported per layer''.

\paragraph{Precision.} float32 arithmetic (relative $10^{-7}$ per operation) propagates the covariance through
sixteen layers of $1024$-term sums to a relative $10^{-6}$, far below the $10^{-4}$ rms of a $10^{-8}$-MSE
estimator; but the quantities that are small differences of $O(1)$ numbers (the Euler--Stein defect
$e-\Delta e$, differences of nearly equal slices) must be formed in float64 or arranged to avoid cancellation.

\section{Measurements on the official networks}
\label{sec:measure}
All numbers are for the official network $0$ ($n=1024$, $L=16$, reference means from $10^9$ Monte Carlo inputs)
unless stated; network $1$ is used as the held-out check where available.

\subsection{The error ladder of the clean chain}
\begin{center}\small
\begin{tabular}{lcc}\toprule
estimator & final-layer MSE & what is added\\\midrule
Gaussian closure (\texttt{gauss\_chain}) & $4.06\cdot10^{-6}$ & --\\
$\kap_3$ CP legs, stars only, no hub & $3.76\cdot10^{-6}$ & births $w$, arms $\diag(g)Ce_m$\\
$+$ hub: $D21\to$ covariance \eqref{eq:cov-transport}, GC arm & $8.37\cdot10^{-7}$ & transported slices\\
$+$ exact birth slice and diagonal (residual leg $R$), exact gates & $8.52\cdot10^{-7}$ & Theorem~\ref{thm:transport}, first term\\
$+$ transposed-slice feed, $\kap_4\!\to\!\kap_3$ feed, $(2,2)$ channel, exact diagonal $\kap_3$ & $5.74\cdot10^{-7}$ & Theorem~\ref{thm:transport}, second and third terms\\
Gaussian closure $+$ Euler--Stein defect, $e-a\delta(0)$, $a=0.535$ & $4.35\cdot10^{-7}$ & Corollary~\ref{cor:eldan}, exact $\delta$ ($2n+1$ chains)\\
Gaussian closure $+$ midpoint $(e+\Delta e)/2$ (no fitted constant) & $4.50\cdot10^{-7}$ & same\\\midrule
organisers' chain, $k_{\max}=3$ simple, numpy port (\texttt{kprop3}), radial $\kap_4$ off & $5.63\cdot10^{-7}$ & exact first-order $\kap_3$ transport, 36 diagrams\\
same, radial $\kap_4$ on (one scalar per layer) & $3.53\cdot10^{-8}$ & the central variable (Lemma~\ref{lem:scale})\\\midrule
open-source third-cumulant chain (504aldo), float32, $33$ fitted scalars, $0.25B$ & $2.13\cdot10^{-8}$ & \\
exact first-order $K=3$ reference of the organisers (no $\kap_4$, no fits) & $4.3\cdot10^{-8}$ & \\
leaders (raw) & $1.1$--$1.5\cdot10^{-8}$ & \\\bottomrule
\end{tabular}
\end{center}
The decisive line is the pair in the middle: the organisers' exact first-order chain, ported to numpy and matched
to the torch reference bit for bit, reaches $5.6\cdot10^{-7}$ on network $0$ when its fourth cumulant is dropped,
the accuracy of our own chain, and $3.5\cdot10^{-8}$ when the single radial scalar
$c=\tfrac{3}{n(n+2)}(\sum_{a\ne b}\kap_4(h_a,h_a,h_b,h_b)+\sum_a\kap_4(h_a))$ per layer is kept. One number per layer,
the collective scale of Lemma~\ref{lem:scale}, is worth a factor $16$ in MSE; our chain's transported $(2,2)$ slice,
which carries $n^2$ numbers and is harmful off the diagonal, was worth nothing. The reference's own naive count is
$3.4\cdot10^{12}$ FLOPs ($1.5B$ in float32), $93\%$ of it in the transport and readout of the factored $\kap_3$ whose rank
grows by $2n$ per layer; the cost problem is the memory of the third cumulant, not the fourth.

Two readings. First, the Euler--Stein defect of the \emph{Gaussian} closure, with no cumulant machinery at all,
reaches the accuracy of our third-cumulant chain: on network $0$ the exact $\delta(0)$ (coordinate Laplacian,
$2049$ chain evaluations on Modal) has correlation $0.931$ with the error, explains $89\%$ of its mean square,
and the fitted $a=0.535$ is close to the parameter-free midpoint $a=\tfrac12$. Theorem~\ref{thm:defect} is
exact; the $11\%$ not explained is the variation of $\delta(y)$ over the Eldan path. Second, the organisers'
exact first-order $K=3$ reference is $13\times$ better than our chain and carries no fitted constant, so the
gap is in the first-order transport, not in calibration or in $\kap_4$.

\subsection{Where the chain's error is injected}
With $e_l=\mathrm{out}_l-u_l$ and the injected part $e_l-\Phi_l\!\cdot\!(W_le_{l-1})$ (the error not explained by
the linear propagation of the previous layer's error through the gates), \texttt{scripts/injected.py} gives rms
injected errors of $2.0,1.1,1.6,2.1,2.3,2.5,2.9,2.7,2.8,2.9,2.8,2.7,2.8,2.9,3.0\cdot10^{-4}$ at layers $1$--$15$
against a total rms error growing from $2.0\cdot10^{-4}$ to $7.6\cdot10^{-4}$: the chain injects a nearly constant
$3\cdot10^{-4}$ per layer from layer $4$ on, and the injected error has no coherent component along the mean
direction (share $<10^{-2}$ beyond layer $2$). This is the signature of a \emph{per-neuron}, incoherent deficit
in the transported slices, consistent with Section~\ref{sec:central}: the transported $D21$ has correlation
$0.965$ with Monte Carlo, i.e.\ an incoherent relative error of $\sqrt{1-0.965^2}\approx26\%$ on top of the
$8$--$10\%$ coherent deficit, and by the empirical law of \cite{Wu}'s reimplementation (extra MSE
$\approx4\cdot10^{-6}\varepsilon^2$ for a relative rms slice error $\varepsilon$) a slice error of $35\%$ accounts for
the whole $5.7\cdot10^{-7}$.

\subsection{What the open-source chain does that ours does not}
A reading of the open-source chain's derivation notes against \texttt{k3chain3.py} (delegated, cross-checked on
the formulas) locates the difference in exactly the first-order transport: (i) its slice update carries every
elementwise diagram of first order in the cumulants, i.e.\ the $\rho$-dependence of the gates
($D21_{ac}\,C_{ac}^k\times$Wick coefficients) and the $\kap_3(z_a)$-driven terms, where ours keeps the gates at
$\rho=0$; each such term is $O(\rho)\approx3\%$ per layer and compounds over fifteen transports; (ii) its fourth
cumulant is the exact diagonal plus an isotropic off-diagonal $\lambda_lC^{\rm off}$ with one fitted $\lambda_l$ per
layer, and it \emph{feeds} $\kap_4$ into the fully off-diagonal $\kap_3$ (the largest single term beyond the base
chain, $2\times$ in MSE), whereas our transported $(2,2)$ slice is measured by them to be harmful off the
diagonal ($1.47\times$ worse than isotropic); (iii) the feedback coefficients of the GC arm are
$(\tfrac12,\tfrac12)$ by the first-order Edgeworth count against our $(1,\tfrac12)$; on network $0$ changing
ours to $\tfrac12$ moves the MSE from $5.74$ to $6.04\cdot10^{-7}$, so this is not the bottleneck. The
$\rho$-exact gates of (i) are the closed forms of Section~\ref{sec:births}: by Lemma~\ref{lem:shift} the
coefficient of $D21_{ac}$ in $\kap_3(h_a,h_a,h_c)$ is $G_{ac}-\mu_aw_a\Phi\big((\alpha_c-\rho_{ac}\alpha_a)/\sqrt{1-\rho_{ac}^2}\big)$
with $G_{ac}=\E[\theta_a\theta_c]$ the gate covariance \eqref{eq:gate}, which reduces to $t_a\Phi_c$ at $\rho=0$,
and likewise for the transposed slice and the $(2,2)$ feed with conditional means and densities at $z_c=0$.
Every coefficient of the first-order transport is a Hermite series obtained by index shifts, which is how a
system built on this note generates its term table (Section~\ref{sec:dictionary}).

\section{Dictionary and open problems}
\label{sec:dictionary}
\subsection{Dictionary}
\begin{center}\small
\begin{tabular}{p{0.42\textwidth}p{0.52\textwidth}}\toprule
stage 8 (commutant, modular tilts, leaf space) & stage 9 (\textsc{MLP} mean estimation)\\\midrule
state $\nu$ on the algebra; affine pairing $\nu(F)$ & input law; the exact answer $u(\nu)=\nu(h_L)$\\
localization process (positive martingale in the commutant) & Eldan tilt of the input, directional tilt, pinning; on the neuron side, the layer-by-layer closure\\
driving covariance $C_t$; trickle-down term $KC_tK$ & which errors a scheme sees: the second variation $E''$ against $d\langle\nu\rangle$ (Theorem~\ref{thm:defect})\\
modular tilt of a Gaussian state & shift of the mean; its action on the chaos expansion is the index shift (Lemma~\ref{lem:shift})\\
central (trivial) part of the localization & the collective activity mode $s=\E z^2$, a scale mixture (Lemma~\ref{lem:scale}); the transverse measure is its law\\
Cartan part & the bulk of the covariance, harmless at percent accuracy (Proposition~\ref{prop:sens})\\
Kikuchi--Toeplitz quantization & the heat semigroup in the input mean (Price's theorem); the exact answer is the Berezin transform of the network function, the chain a Toeplitz symbol that is not heat-harmonic, and the anomaly is the Euler--Stein defect\\
capacity of a localization & Proposition~\ref{prop:capacity}: input tilts have signal-to-noise $\sqrt k\,|\delta|/\|A\|_F\approx10^{-2}\sqrt k$\\
births / creation operators & the kink births $w_a=\varphi(\alpha_a)/\sigma_a$, second derivatives of the tilt response (\eqref{eq:birth})\\
\bottomrule\end{tabular}
\end{center}

\subsection{What the theory prescribes for a system, and what is open}
The decisive constraints are the pair floor of Proposition~\ref{prop:floor} (four $n\times n$ objects per layer
besides the covariance, at the $0.1B$ floor) and the error accounting of Section~\ref{sec:central} (the $(2,1)$
transport carries the error; the $(2,2)$ channel is worth $10^{-5}$; the $(3,1)$ term has the wrong sign at first
order). A system built on this note therefore spends its budget on \emph{one} object, the $(2,1)$ slice, transported
with the $\rho$-exact gate, in float32 with the symmetric engine, and corrects the rest by calibration fitted on the
public networks, in this order of expected value:
\begin{enumerate}[nosep]
\item \textbf{$\rho$-exact slice gate.} Replace the gate $g_a-\mu_aw_a$ of \eqref{eq:slice-transport} by its
value at the true correlation $\rho_{ac}$ (one more Hermite series in $\rho$, $O(Kn^2)$, no new $n^3$ object). The
$8$--$10\%$ slice deficit is of the size of the $O(\rho)$ correction along the collective mode.
\item \textbf{Euler--Stein merge.} $e-a\,\delta(0)$ with $a$ fitted once: a factor $3$ (\texttt{k3v3}) to $15$
(Gaussian closure) at $n=256$; at $n=1024$ the exact $\delta$ of the Gaussian closure is measured in
Section~\ref{sec:measure}. Within budget only through the closable part of the second-order response of
Theorem~\ref{thm:gram}, i.e.\ the covariances $\E[\delta m\delta m^\top]$, $\E[\delta m_a\delta S_{ab}]$,
$\E[\delta m_a\delta S_{bb}]$, $\E[\delta S_{aa}\delta S_{bb}]$, $\E[\delta S_{ab}\delta S_{aa}]$, which are the
same five objects a $\kap_3$ chain carries under other names; whether the closable part correlates with the exact
$\delta$ as well as $\delta$ correlates with the error is the open experiment of Section~\ref{sec:measure}.
\item \textbf{Per-layer slice scale.} A fitted multiplier $1+\gamma_l$ on the transported slices (the chain's
\texttt{d3scale} generalised per layer): one number per layer, fitted on the public networks, directly targets the
measured deficit.
\item \textbf{Noncommutative version.} The gates $\theta(z_a)$ of one layer generate a commutative algebra; across
layers the gate algebras do not commute with the weights, and the chain's closure is a product state on the gate
algebra of layer $l$ conditioned on the surrogate of layer $l-1$. The stage-8 KMS condition for the modular flow of
the Gaussian state, restricted to the gate algebra, is the heat equation of Section~\ref{sec:births}; its failure
across layers is the defect. Writing the chain as a state on the free product of the gate algebras, with the
trickle-down of \cite{AKV} as its Cartan expansion, is the theoretical continuation.
\end{enumerate}

\appendix
\section{Verification}
\texttt{scripts/verify\_stage9.py} checks, in order: (1) Lemma~\ref{lem:shift} to $2\cdot10^{-8}$ for $k\le10$;
(2) Price's theorem and Stein's lemma for the bivariate \textsc{ReLU} moment,
$\partial_{m_a}M_{ab}=\E[\theta_a\relu z_b]$, $\partial_{S_{ab}}M_{ab}=\E[\theta_a\theta_b]=\partial_{m_a}\partial_{m_b}M_{ab}$,
against $4\cdot10^6$ Monte Carlo samples; (3) Theorem~\ref{thm:gram} on a network with $n=6$, $L=3$: the gradients
of the first-layer moments in closed form against finite differences ($1.6\cdot10^{-7}$ on a scale of $2$), and
$\delta(0)$ by finite differences in the input mean against $-\langle\nabla^2\mathcal G,\Gamma\rangle$ with the
Hessian of the downstream chain by finite differences in the $27$ moment coordinates: agreement to
$7\cdot10^{-7}$ on $|\delta|_{\rm rms}=0.28$; (4) Lemma~\ref{lem:scale} by Gauss--Hermite quadrature in the
scale, $v=10^{-2}$: the closure error matches $-\tfrac v2\sigma(\alpha^2-1)\varphi$ and the $\kap_4$ readout cancels
it to the $O(v^2)$ level for $\alpha\in\{-1,0,0.7,1.5\}$.
\texttt{scripts/defect\_small.py}, \texttt{lap\_exact.py}, \texttt{lapx\_merge.py} produce the numbers of
Sections~\ref{sec:capacity} and~\ref{sec:measure}; \texttt{diag\_eig\_v3.py}, \texttt{diag\_onestep.py},
\texttt{injected.py} those of Section~\ref{sec:central}.

\bibliographystyle{plain}
\begin{thebibliography}{99}
\bibitem{stage8} Stage 8 note: Localization in the commutant: modular tilts, Kikuchi--Toeplitz quantization and transverse measures (9 October 2026).
\bibitem{Wu} W.~Wu, V.~Lecomte, M.~Winer, G.~Robinson, J.~Hilton, P.~Christiano, Estimating the expected output of wide random MLPs more efficiently than sampling, arXiv:2605.05179.
\bibitem{AKV} N.~Anari, F.~Koehler, T.-D.~Vuong, Trickle-down in localization schemes and applications, arXiv:2407.16104.
\bibitem{CE} Y.~Chen, R.~Eldan, Localization schemes: a framework for proving mixing bounds for Markov chains, FOCS 2022.
\bibitem{Hanin} B.~Hanin, Random fully connected neural networks as perturbatively solvable hierarchies, arXiv:2204.01058.
\bibitem{Connes} A.~Connes, Sur la th\'eorie non commutative de l'int\'egration, LNM 725 (1979).
\end{thebibliography}
\end{document}
